\chapter{12.10-2 等腰三角形的性质与判定}


\section*{二、填空题（共6小题）}

\blanklist[9]

\item 如图，在 $\triangle ABC$ 中，$\angle BAC = 150^\circ$，$AD \perp BC$ 于 $D$，且 $AB + BD = DC$，那么 $\angle C =$ \_\_\_\_\_\_。

\rightquestionimage[width=0.35\textwidth]{0.35\textwidth}{12.10-2-9.png}

\item 如图，$\triangle ABC$ 中，$AB = AC$，$BD$ 平分 $\angle ABC$ 交 $AC$ 于 $G$，$DM \parallel BC$ 交 $\angle ABC$ 的外角平分线于 $M$，交 $AB$、$AC$ 于 $F$、$E$，下列结论：

① $MB \perp BD$；② $FD = EC$；③ $EC = EF + DG$；④ $CE = \frac{1}{2}MD$。

其中正确的有\_\_\_\_\_\_（填序号）。

\rightquestionimage[width=0.35\textwidth]{0.35\textwidth}{12.10-2-10.png}

\item 如图：已知 $\angle BAD = \angle DAC = 90^\circ$，$AD \perp AE$，且 $AB + AC = BE$，则 $\angle B =$ \_\_\_\_\_\_。

\rightquestionimage[width=0.35\textwidth]{0.35\textwidth}{12.10-2-11.png}

\item 等腰三角形一腰上的高等于该三角形某一条边的长度的一半，则其顶角可能为\_\_\_\_\_\_。

\item 如图，点 $B$ 是线段 $AC$ 的中点，过点 $C$ 的射线 $CE$ 与 $AC$ 成 $60^\circ$ 的角，点 $P$ 为射线 $CE$ 上一动点，给出以下四个结论：

① 当 $AP \perp CE$，垂足为 $P$ 时，$\angle APB = 30^\circ$；

② 当 $CP = AC$ 时，$\angle APB = 30^\circ$；

③ 在射线 $CE$ 上，使 $\triangle APC$ 为直角三角形的点 $P$ 只有 1 个；

④ 在射线 $CE$ 上，使 $\triangle APC$ 为等腰三角形的点 $P$ 只有 1 个；

其中正确结论的序号是\_\_\_\_\_\_。

\rightquestionimage[width=0.35\textwidth]{0.35\textwidth}{12.10-2-13.png}

\item 如图，$\triangle ABC$ 中，$AB = BC = AC = 12\text{cm}$，现有两点 $M$、$N$ 分别从点 $A$、$B$ 同时出发，沿三角形的边运动，已知点 $M$ 的速度为 $1\text{cm/s}$，点 $N$ 的速度为 $2\text{cm/s}$。当点 $N$ 第一次到达 $B$ 点时，$M$、$N$ 同时停止运动。


(1) 当点 $M$、$N$ 运动\_\_\_\_\_\_秒时，$M$、$N$ 两点重合；

(2) 当点 $M$、$N$ 运动\_\_\_\_\_\_秒后，$M$、$N$ 与 $\triangle ABC$ 中的某个顶点可得到等腰三角形。

\rightquestionimage[width=0.35\textwidth]{0.35\textwidth}{12.10-2-14.png}

\end{enumerate}

\section*{三、解答题（共3小题）}

\solutionlist[15]

\item 已知：如图，$AD$ 是 $\triangle BAC$ 的角平分线，$BE \perp AD$ 交 $AD$ 的延长线于点 $E$，$EF \parallel AC$ 交 $AB$ 于点 $F$，$AE = 8$，$EF = 5$。求 $BE$ 的长。

\rightquestionimage[width=0.35\textwidth]{0.35\textwidth}{12.10-2-15.png}

\newpage

\item 在 $\triangle ABC$ 中，$\angle C = 90^\circ$，$AC = BC = 2$，将一块三角板的直角顶点放在斜边 $AB$ 的中点 $P$ 处，将此三角板绕点 $P$ 旋转，三角板的两直角边分别交射线 $AC$、$CB$ 与点 $D$、点 $E$，图①、②、③是旋转得到的三种图形。

(1) 观察线段 $PD$ 和 $PE$ 之间的有怎样的关系，并以图②为例，加以说明；

(2) $\triangle PBE$ 是否构成等腰三角形？若能，指出所有的情况（即求出 $\triangle PBE$ 为等腰三角形时 $CE$ 的长）；若不能请说明理由。

\centerquestionimage[width=0.70\textwidth]{0.70\textwidth}{12.10-2-16.png}

\vspace{5cm}

\item 如图，在 $\triangle ABC$ 中，$AB = AC$，$P$ 为 $\triangle ABC$ 内一点，且 $\angle BAP = 70^\circ$，$\angle ABP = 40^\circ$。


(1) 求证：$\triangle ABP$ 是等腰三角形；

(2) 连接 $PC$，当 $\angle PCB = 30^\circ$ 时，求 $\angle PBC$ 的度数。

\rightquestionimage[width=0.40\textwidth]{0.40\textwidth}{12.10-2-17.png}

\end{enumerate}
